\documentclass[10pt,a4paper]{amsart}
\usepackage[margin=19mm]{geometry}
\usepackage{amsmath,amssymb,mathtools}
\usepackage[scr=rsfs,cal=euler]{mathalfa}
\usepackage{xfrac}
\usepackage[unicode,hidelinks,bookmarks=false]{hyperref}
\newcommand{\ie}{i.e.\ }
\newcommand{\cf}{cf.\ }
\newcommand{\resp}{respectively\ }
\newcommand{\Subsection}{\subsection}
\providecommand{\nobreakdash}{\nobreak\hbox{-}}
\setlength{\emergencystretch}{2em}
\multlinegap0pt
\usepackage{tikz}
\usepackage{xfrac}
\usepackage{amssymb}
\usepackage{mathtools}
\usepackage[shortlabels]{enumitem}
\usepackage{float}
\newtheorem{lemma}{Lemma}
\newtheorem{corollary}{Corollary}
\newcommand\A{\mathsf{A}}
\newcommand\C {C}
\newcommand{\E}{E}
\newcommand{\ICK}{\mathsf{ICK}}
\newcommand{\K}{K}
\newcommand{\LICK}{\mathcal{L}_{\ICK}}
\newcommand\Prop {\mathsf{Prop}}
\newcommand{\Rel}{\mathrel{R}}
\newcommand{\W}{\mathcal{W}}
\newcommand{\agi}{a}
\newcommand{\fw}{\mathcal{M}}
\newcommand{\hick}{\mathrm{ICK}}
\newcommand{\hickp}{\mathrm{P}}
\newcommand{\hicks}{\mathrm{ICKS4}}
\newcommand{\hickss}{\mathrm{ICKS5}}
\newcommand{\hickt}{\mathrm{ICKT}}
\title{Intuitionistic Common Knowledge}
\author{Lukas Zenger}
\date{Source: arXiv:2605.00523v1 (CC BY 4.0)}
\begin{document}
\maketitle

\noindent\emph{Source and licence.} Intuitionistic Common Knowledge. Version arXiv:2605.00523v1 (CC BY 4.0), distributed by arXiv under the Creative Commons Attribution 4.0 International licence, http://creativecommons.org/licenses/by/4.0/, as recorded in the arXiv licence metadata for this exact version and shown on the abstract page https://arxiv.org/abs/2605.00523v1. Full licence notice: \texttt{licenses/arxiv-2605.00523-CC-BY-4.0.txt}.\par\smallskip

\noindent\textit{Editorial scope.} Counted unit: the article's lemma lemma (source0.tex lines 856--861). The article's complete original proof of it is delivered with it (source0.tex lines 862--996, one \texttt{proof} environment of 135 lines). No mathematics is written, repaired or re-derived: the statement and the proof are the article's own lines, in the article's own notation, and no retained line is substituted or re-flowed.  Retained uncounted as the source-local context the proof needs: lines 495--502; lines 508--510; lines 552--557; lines 585--593; lines 749--758; lines 844--851.  Nothing was omitted from any retained span, so the package carries no omission record.  No retained line contains a citation, so the wrapper carries no bibliography and no bibliography entry.  No retained line contains a figure environment, an image-inclusion command or drawing code, so no image is delivered and the package declares no asset dependency.  Editorial formatting only, in addition: an AMS \texttt{amsart} preamble that restates the article's own declarations and macros used by the retained lines, namely the environments \texttt{lemma}, \texttt{corollary} on separate counters, together with the standard packages those lines need.  The retained environments print in this document as Lemma 1, Corollary 1 in order of appearance, whereas the article prints the same items with the numbering of its own full document.  The complete native source of the article as distributed in the arXiv e-print for this exact version is bundled unchanged as \texttt{sources/199570/source0.tex}.
\begin{lemma}\label{l: derivable formulae}
    Let $\varphi, \psi \in \LICK$ and $\hickp \in \{\hick, \hickt, \hicks, \hickss\}$.
    \begin{enumerate}
    \item\label{l: derivable formulae item 1} If $\vdash_\hickp \varphi$, then $\vdash_\hickp \gamma \rightarrow \varphi$ for any $\gamma \in \LICK$.
    \item\label{l: derivable formulae item 2} If $\vdash_{\hickp} \varphi \rightarrow \psi$ and $\vdash_{\hickp} \psi \rightarrow \chi$, then $\vdash_{\hickp} \varphi \rightarrow \chi$.
    \item\label{l: derivable formulae item 3} $\vdash_\hickp (\varphi \rightarrow \psi) \rightarrow (\neg \psi \rightarrow \neg \varphi)$ (contraposition).\footnote{Note that the other direction, i.e. $(\neg \psi \rightarrow \neg \varphi) \rightarrow (\varphi \rightarrow \psi)$ is \emph{not} valid intuitionistically!}
    \end{enumerate}
\end{lemma}

\begin{lemma}[Assumption weakening] \label{l: assumption weakening}
    Let $\hickp \in \{\hick, \hickt, \hicks, \hickss\}$. If $\Gamma \vdash_\mathrm{P} \varphi$ and $\Gamma \subseteq \Delta$, then $\Delta \vdash_\mathrm{P} \varphi$.
\end{lemma}

\begin{corollary}\label{c: deduction theorem}
    For any finite set of formulae $\Gamma \cup \{\varphi\}$ the following holds:
    \begin{center}
        $\Gamma \vdash_\mathrm{P} \varphi$ if and only if $\vdash_\mathrm{P} \bigwedge \Gamma \rightarrow \varphi$.
    \end{center}
\end{corollary}

\begin{lemma}\label{l: derivable formulae 2}
   Let $\varphi, \psi \in \LICK$, $\agi \in \A$ and $\hickp \in \{\hick, \hickt, \hicks, \hickss\}$.
   \begin{enumerate}
       \item\label{l: derivable formulae 2 item 1} $\vdash_{\hickp} \K_\agi (\varphi \wedge \psi) \rightarrow (\K_\agi \varphi \wedge \K_\agi \psi)$ and $\vdash_{\hickp} (\K_\agi \varphi \wedge \K_\agi \psi) \rightarrow \K_\agi (\varphi \wedge \psi)$.
    \item\label{l: derivable formulae 2 item 2} $\vdash_{\hickp} (\K_\agi \varphi \vee \K_\agi \psi) \rightarrow \K_\agi (\varphi \vee \psi)$.
    \item\label{l: derivable formulae 2 item 3} If $\vdash_{\hickp} \varphi$, then $\vdash_{\hickp} \C \varphi$ (necessitation for $\C$).
    \item\label{l: derivable formulae 2 item 4} $\vdash_{\hickss} K_\agi \varphi \vee \neg K_\agi \varphi$.
   \end{enumerate}
\end{lemma}

\begin{lemma}\label{l: properties of prime theories}
    Let $\Gamma$ be a $\Sigma$-prime theory (relative to $\mathrm{P}$).
    \begin{enumerate}
        \item\label{l: properties of prime theories, Item 1} If $\varphi \wedge \psi \in \Sigma$, then $\varphi \wedge \psi \in \Gamma$ if and only if $\varphi \in \Gamma$ and $\psi \in \Gamma$.
        \item\label{l: properties of prime theories, Item 2} If $\varphi \vee \psi \in \Sigma$, then $\varphi \vee \psi \in \Gamma$ if and only if $\varphi \in \Gamma$ or $\psi \in \Gamma$.
        \item\label{l: properties of prime theories, Item 3}  $\C \varphi \in \Gamma$ if and only if $\varphi \in \Gamma$ and $\{\K_\agi \C \varphi \mid \agi \in \A\} \subseteq  \Gamma$.
        \item\label{l: properties of prime theories, Item 4} If $\mathsf{T_\agi} \in \mathrm{P}$, then $\K_\agi \varphi \in \Gamma$ implies $\varphi \in \Gamma$.
        %\item\label{l: properties of prime theories, Item 5}
    \end{enumerate}
\end{lemma}

\begin{lemma}\label{l: properties characteristic formula IM}
    Let $\Gamma \in \W^{(\Sigma, \hickp)}$.  The following hold.
    \begin{enumerate}
        \item $\Gamma \vdash_\mathrm{P} \chi(\Gamma)$.
        \item For any $\Delta \in (R^{(\Sigma,\hickp)})^*(\Gamma)$, $\Delta \vdash_\mathrm{P} \gamma(\Gamma)$.
        \item If $\psi \in \LICK$ such that $\psi \in \Delta$ for every $\Delta \in R^{(\Sigma, \hickp)^*}(\Gamma)$, then $\vdash_\mathrm{P} \gamma(\Gamma) \rightarrow \psi$.
    \end{enumerate}
\end{lemma}

\begin{lemma}[Truth Lemma] \label{l: truth lemma}
Let $\hickp \in \{\hick, \hickt, \hicks, \hickss\}$, $\Sigma$ a finite, non-empty and closed set of formulae (negation closed if $\hickp = \hickss$), and $\fw^{(\Sigma, \hickp)}$ the canonical model relative to $\Sigma$ and $\hickp$. Then for any $\Gamma \in \W^{(\Sigma, \hickp)}$ and any formula $\varphi \in \Sigma$ the following holds.
\begin{center}
    $\varphi \in \Gamma$ if and only if $\fw^{(\Sigma, \hickp)}, \Gamma \models \varphi$.
\end{center}
\end{lemma}

\begin{proof}
    Let $\hickp \in \{\hick, \hickt, \hicks, \hickss\}$. We proceed by induction on the structure of $\varphi$. The case where $\varphi = \bot$ is trivial, and the case where $\varphi = p$ for $p \in \Prop$ follows immediately from the definition. For the induction step, the cases where $\varphi = \psi \wedge \chi$ or $\varphi = \psi \vee \chi$ follow immediately from Lemma \ref{l: properties of prime theories}, the induction hypothesis and $\Sigma$ being closed. \smallskip
    
    \noindent \textsc{Case for $\rightarrow$.} Suppose $\varphi = \psi \rightarrow \chi$. For the left-to-right direction suppose $\psi \rightarrow \chi \in \Gamma$. Let $\Delta \in \W^{(\Sigma, \hickp)}$ with $\Gamma \subseteq \Delta$ and suppose $\fw^{(\Sigma, \hickp)}, \Delta \models \psi$. By induction hypothesis $\psi \in \Delta$ implying that $\Delta \vdash_\mathrm{P} \psi$. Since $\Delta$ extends $\Gamma$, we have $\psi \rightarrow \chi \in \Delta$ and thus $\Delta \vdash_\mathrm{P} \psi \rightarrow \chi$. Applying $\mathsf{MP}$ yields $\Delta \vdash_\mathrm{P} \chi$. As $\Delta$ is deductively closed with respect to $\Sigma$, $\chi \in \Delta$. The induction hypothesis yields $\fw^{(\Sigma, \hickp)}, \Delta \models \chi$. $\Delta$ was an arbitrary $\Sigma$-prime theory extending $\Gamma$, therefore $\fw^{(\Sigma, \hickp)}, \Gamma \models \psi \rightarrow \chi$. 
    For the right-to-left direction we proceed by contraposition. Suppose $\psi \rightarrow \chi \not \in \Gamma$. Then $\Gamma \not \vdash_\mathrm{P} \psi \rightarrow \chi$. The Deduction Theorem yields $\Gamma, \psi \not \vdash_\mathrm{P} \chi$. By the Lindenbaum Lemma there exists $\Delta \in \W^{(\Sigma, \hickp)}$ with $\Gamma \cup \{\psi\} \subseteq \Delta$ and $\Delta \not \vdash_\mathrm{P} \chi$. Hence, $\psi \in \Delta$ and $\chi \not \in \Delta$. The induction hypothesis yields $\fw^{(\Sigma, \hickp)}, \Delta \models \psi$ and $\fw^{(\Sigma, \hickp)}, \Delta \not \models \chi$. As $\Gamma \leq^{(\Sigma, \hickp)} \Delta$, we have $\fw^{(\Sigma, \hickp)}, \Gamma \not \models \psi \rightarrow \chi$. \smallskip
    
    \noindent \textsc{Case for $\K_\agi$.} Suppose $\varphi = \K_\agi \psi$. We distinguish the following cases depending on $\hickp$. \smallskip
   
    \setlength{\leftskip}{0.5cm} \noindent \textsc{Case for $\hickp \in \{\hick, \hickt\}$}.  For the left-to-right direction suppose $\K_\agi \psi \in \Gamma$. Let $\Delta \in \W^{(\Sigma, \hickp)}$ be any $\Sigma$-prime theory such that $\Gamma \Rel_\agi^{(\Sigma, \hickp)} \Delta$ holds. By definition $\K_\agi^{-1} \Gamma \subseteq \Delta$. Hence, $\psi \in \Delta$. By induction hypothesis $\fw^{(\Sigma, \hickp)}, \Delta \models \psi$. As $\Delta$ was arbitrary, we conclude that $\fw^{(\Sigma, \hickp)}, \Gamma \models \K_\agi \psi$. For the right-to-left direction suppose $\fw^{(\Sigma, \hickp)}, \Gamma \models \K_\agi \psi$. Note that
    \begin{equation}\label{e: truth lemma1}
        \K_\agi^{-1} \Gamma \vdash_\mathrm{P} \psi
    \end{equation}
     as otherwise, by the Lindenbaum Lemma, there would exist a $\Sigma$-prime theory $\Delta$ with $\K_\agi^{-1} \Gamma \subseteq \Delta$ and $\Delta \not \vdash_\mathrm{P} \psi$. Hence $\psi \not \in \Delta$, and so by induction hypothesis $\fw^{(\Sigma, \hickp)}, \Delta \not \models \psi$, contradicting that $\fw^{(\Sigma, \hickp)}, \Gamma \models \K_\agi \psi$. Therefore (\ref{e: truth lemma1}) holds. Since $\K_\agi^{-1}\Gamma$ is a finite set, from (\ref{e: truth lemma1}) and Corollary \ref{c: deduction theorem} we obtain
    \begin{equation}\label{e: truthlemma2}
        \vdash_\mathrm{P} \bigwedge \K_\agi^{-1} \Gamma \rightarrow \psi.
    \end{equation}
    By $\mathsf{Nec_\agi}$ we obtain $\vdash_\mathrm{P} \K_\agi (\bigwedge \K_\agi^{-1} \Gamma \rightarrow \psi) $ and therefore $\vdash_\mathrm{P} \K_\agi (\bigwedge \K_\agi^{-1} \Gamma) \rightarrow \K_\agi \psi$ by the axiom $\mathsf{K_\agi}$ and $\mathsf{MP}$. We claim that
    \begin{equation}\label{e: truthlemma3}
         \K_\agi \K_\agi^{-1} \Gamma \vdash_\mathrm{P} \K_\agi (\bigwedge \K_\agi^{-1} \Gamma).
    \end{equation}
    If true, then applying $\mathsf{MP}$ to $\vdash_\mathrm{P} \K_\agi (\bigwedge \K_\agi^{-1} \Gamma) \rightarrow \K_\agi \psi$ and (\ref{e: truthlemma3}) gives $\K_\agi \K_\agi^{-1} \Gamma \vdash \K_\agi \psi$. Lemma \ref{l: assumption weakening} then yields $\Gamma \vdash \K_\agi \psi$ implying that $\K_\agi \psi \in \Gamma$. We prove (\ref{e: truthlemma3}) by induction on the size of $\K_\agi^{-1} \Gamma$. For $\lvert \K_\agi^{-1} \Gamma \rvert = 0$, the statement reduces to $\emptyset \vdash \K_\agi \top$ which is trivially true. For $\lvert \K_\agi^{-1} \Gamma \rvert = k+1$ let $\K_\agi^{-1}\Gamma_0 = \{ \varphi_1, \ldots, \varphi_k\}$ and let $\K_\agi^{-1}\Gamma = \K_\agi^{-1} \Gamma_0 \cup \{ \varphi_{k+1}\}$. By induction hypothesis
    \begin{equation*}
        \K_\agi \K_\agi^{-1} \Gamma_0 \vdash_\mathrm{P} \K_\agi (\bigwedge \K_\agi^{-1} \Gamma_0)
    \end{equation*}
    Hence also $\K_\agi \K_\agi^{-1} \Gamma \vdash_\mathrm{P} \K_\agi (\bigwedge \K_\agi^{-1} \Gamma_0)$. Since $\K_\agi \K_\agi^{-1} \Gamma \vdash_\mathrm{P} \K_\agi \varphi_{k+1}$ we also obtain
    \begin{equation*}
        \K_\agi \K_\agi^{-1} \Gamma \vdash_\mathrm{P} \K_\agi (\bigwedge \K_\agi^{-1} \Gamma_0) \wedge \K_\agi \varphi_{k+1}.
    \end{equation*}
     Lemma \ref{l: derivable formulae}, Item \ref{l: derivable formulae item 2}. yields
     \begin{equation*}
         \K_\agi \K_\agi^{-1} \Gamma \vdash_\mathrm{P} \K_\agi (\bigwedge \K_\agi^{-1} \Gamma)
     \end{equation*}
     which finishes the proof of the claim. \smallskip

     \noindent \textsc{Case for $\hickp = \hicks$.} For the left-to-right direction suppose $\K_\agi \psi \in \Gamma$. Let $\Delta \in \W^{(\Sigma, \hicks)}$ be any $\Sigma$-prime theory such that $\Gamma \Rel_\agi^{(\Sigma, \hicks)} \Delta$ holds. By definition $\K_\agi \K_\agi^{-1} \Gamma \subseteq \Delta$. Hence $\K_\agi \psi \in \Delta$. Since $\Delta \vdash_\hicks \K_\agi \psi \rightarrow \psi$, applying $\mathsf{MP}$ gives $\Delta \vdash_\hicks \psi$ and so $\psi \in \Delta$. The induction hypothesis yields $\fw^{(\Sigma, \hicks)}, \Delta \models \psi$. As $\Delta$ was arbitrary $\fw^{(\Sigma, \hicks)}, \Gamma \models \K_\agi \psi$. For the right-to-left direction suppose $\fw^{(\Sigma, \hicks)}, \Gamma \models \K_\agi \psi$. Note that
    \begin{equation}\label{e: completeness S4 K1}
       \K_\agi \K_\agi^{-1} \Gamma \vdash_\hicks \psi
    \end{equation}
    as otherwise there would exist, by the Lindenbaum Lemma, a prime theory $\Delta \supseteq \K_\agi \K_\agi^{-1} \Gamma$ with $\Delta \not \vdash_\hicks \psi$. Hence $\psi \not \in \Delta$ and so by induction hypothesis $\fw^{(\Sigma, \hicks)}, \Delta \not \models \psi$. By construction $\Gamma \Rel_\agi^{(\Sigma, \hicks)} \Delta$, implying that $\fw^{(\Sigma, \hicks)}, \Gamma \not \models \K_\agi \psi$; contradiction. Hence (\ref{e: completeness S4 K1}) holds. By Corollary \ref{c: deduction theorem} we have
    \begin{equation}\label{e: completeness S4 K2}
        \vdash_\hicks \bigwedge \K_\agi \K_\agi^{-1} \Gamma \rightarrow \psi.
    \end{equation}
    By $\mathsf{Nec_\agi}$, $\mathsf{K_\agi}$ and $\mathsf{MP}$ we obtain
    \begin{equation}\label{e: completeness S4 K3}
        \vdash_\hicks \K_\agi \bigwedge \K_\agi \K_\agi^{-1} \Gamma \rightarrow \K_\agi \psi.
    \end{equation}
   We claim that $\vdash_\hicks \bigwedge \K_\agi \K_\agi \K_\agi^{-1} \Gamma \rightarrow \K_\agi \bigwedge \K_\agi \K_\agi^{-1} \Gamma$ and we prove the claim by induction on $\lvert \K_\agi^{-1} \Gamma \rvert$. For ${\lvert \K_\agi^{-1} \Gamma \rvert = 0}$ the statement reduces to $\vdash_\hicks  \top \rightarrow \K_\agi \top$, which is easy to prove. For ${\lvert \K_\agi^{-1} \Gamma \rvert = k+1}$, let $\K_\agi^{-1} \Gamma = \K_\agi^{-1}\Gamma_0 \cup \{\varphi_{k+1}\}$ where $ \K_\agi^{-1} \Gamma_0= \{\varphi_1, \ldots,\varphi_k\}$. By induction hypothesis
   \begin{equation}\label{e: completeness S4 K4}
       \vdash_\hicks \bigwedge \K_\agi \K_\agi \K_\agi^{-1}\Gamma_0 \rightarrow \K_\agi \bigwedge \K_\agi \K_\agi^{-1} \Gamma_0.
   \end{equation}
   Since $(p_1 \rightarrow q_1) \rightarrow ((p_2 \rightarrow q_2) \rightarrow ((p_1 \wedge p_2) \rightarrow (q_1 \wedge q_2)))$ is an intuitionistic tautology and $\vdash_\hicks \K_\agi \K_\agi \varphi_{k+1} \rightarrow \K_\agi \K_\agi \varphi_{k+1}$, we obtain
   \begin{equation}\label{e: completeness S4 K5}
       \vdash_\hicks \bigwedge \K_\agi \K_\agi \K_\agi^{-1}\Gamma \rightarrow ((\K_\agi \bigwedge \K_\agi \K_\agi^{-1} \Gamma_0) \wedge \K_\agi \K_\agi \varphi_{k+1})
   \end{equation}
 By Lemma \ref{l: derivable formulae 2}, Item \ref{l: derivable formulae 2 item 1} we have
 \begin{equation}\label{e: completeness S4 K6}
     \vdash_\hicks ((\K_\agi \bigwedge \K_\agi \K_\agi^{-1} \Gamma_0) \wedge \K_\agi \K_\agi \varphi_{k+1}) \rightarrow \K_\agi \bigwedge \K_\agi \K_\agi^{-1} \Gamma.
 \end{equation}
 Hence we we obtain $\vdash_\hicks \bigwedge \K_\agi \K_\agi \K_\agi^{-1} \Gamma \rightarrow \K_\agi \bigwedge \K_\agi \K_\agi^{-1} \Gamma$ by Lemma \ref{l: derivable formulae}, Item \ref{l: derivable formulae item 2}  as claimed. Now since $\Gamma \vdash_\hicks \K_\agi \gamma$ for each $\K_\agi \gamma \in \Gamma$, axiom $\mathsf{S4_\agi}$ implies that $\Gamma \vdash_\hicks \K_\agi \K_\agi \gamma$ for each $\K_\agi \gamma \in \Gamma$ and so, using standard arguments, $\Gamma \vdash_\hicks \bigwedge \K_\agi \K_\agi \K_\agi^{-1} \Gamma$. Hence $\Gamma \vdash_\hicks \K_\agi \bigwedge \K_\agi \K_\agi^{-1} \Gamma$ and so we obtain $\Gamma \vdash_\hicks \K_\agi \psi$ from (\ref{e: completeness S4 K3}). Therefore $\K_\agi \psi \in \Gamma$. \smallskip

 \noindent \textsc{Case for $\hickss$.} For the direction from left-to-right suppose $\K_\agi \psi \in \Gamma$ and let $\Delta \in \W^{(\Sigma, \hickss)}$ be any world such that $\Gamma \Rel_\agi^{(\Sigma, \hickss)} \Delta$. By definition $\K_\agi \K_\agi^{-1} \Gamma = \K_\agi \K_\agi^{-1} \Delta$, implying that $\K_\agi \psi \in \Delta$. Thus $\psi \in \Delta$ by axiom $\mathsf{T_\agi}$ and so $\fw^{(\Sigma, \hickss)}, \Delta \models \psi$ by induction hypothesis. As $\Delta$ was arbitrary we conclude $\fw^{(\Sigma, \hickss)}, \Gamma \models \K_\agi \psi$. For the right-to-left direction suppose $\fw^{(\Sigma, \hickss)}, \Gamma \models \K_\agi \psi$. We claim that
 \begin{equation}\label{e: completeness S5 K1}
     \K_\agi \K_\agi^{-1} \Gamma, \neg \K_\agi \neg \K_\agi^{-1} \Gamma \vdash_\hickss \psi.
 \end{equation}
 In fact if (\ref{e: completeness S5 K1}) would not hold, then by the Lindenbaum Lemma there would exist a $\K$-maximally consistent prime theory $\Delta$ such that (1) $ \K_\agi \K_\agi^{-1} \Gamma \cup \neg \K_\agi \neg \K_\agi^{-1} \Gamma \subseteq \Delta$ and (2) $\Delta \not \vdash_\hickss \psi$. From (2) follows by the induction hypothesis that $\fw^{(\Sigma, \hickss)} \not \models \psi$. From (1) follows that $\Gamma \Rel_\agi^{(\Sigma, \hickss)} \Delta$: By construction $\K_\agi \K_\agi^{-1} \Gamma \subseteq \Delta$. Suppose $\K_\agi \chi \in \Delta$. Since $\Gamma$ is $\K$-maximally consistent we have either (i) $\K_\agi \chi \in \Gamma$ or (ii) $\neg \K_\agi \chi \in \Gamma$. Note that (ii) implies that $\neg \K_\agi \chi \in \Delta$ which implies that $\Delta \vdash_\hickss \bot$. Hence (i) holds, i.e. $\K_\agi \K_\agi^{-1} \Gamma = \K_\agi \K_\agi^{-1} \Delta$. Thus (1) and (2) imply that $\fw^{(\Sigma, \hickss)}, \Gamma \not \models \K_\agi \psi$; a contradiction. We conclude that (\ref{e: completeness S5 K1}) holds. By Corollary \ref{c: deduction theorem} we have
 \begin{equation}\label{e: completeness S5 K2}
     \vdash_\hickss \bigwedge \Omega \rightarrow \psi
 \end{equation}
 where $\Omega = \K_\agi \K_\agi^{-1} \Gamma \cup \neg \K_\agi \neg \K_\agi^{-1} \Gamma$. By $\mathsf{Nec_\agi}$, $\mathsf{K_\agi}$ and $\mathsf{MP}$ we obtain
 \begin{equation}\label{e: completeness S5 K3}
     \vdash_\hickss \K_\agi \bigwedge \Omega \rightarrow \K_\agi \psi.
 \end{equation}
We can then show that 
 \begin{equation}\label{e: completeness S5 K4}
     \vdash_\hickss \bigwedge \K_\agi \Omega \rightarrow \K_\agi \bigwedge \Omega
 \end{equation}
 by induction on $\lvert \Omega \rvert$. The proof is similar to the corresponding case for $\hicks$ and omitted. Now since $\Gamma \vdash_\hickss \K_\agi \gamma$ for each $\K_\agi \gamma \in \Gamma$, axiom $\mathsf{S4_\agi}$ implies $\Gamma \vdash_\hickss \K_\agi \K_\agi \gamma$ for each $\K_\agi \gamma \in \Gamma$. Similarly since $\Gamma \vdash_\hickss \neg \K_\agi \gamma$ for each $\neg \K_\agi \gamma \in \Gamma$, axiom $\mathsf{S5_\agi}$ implies $\Gamma \vdash_\hickss \K_\agi \neg \K_\agi \gamma$ for each $\neg \K_\agi \gamma \in \Gamma$. Thus $\Gamma \vdash_\hickss \bigwedge \K_\agi \Omega$ and so $\Gamma \vdash_\hickss \K_\agi \bigwedge \Omega$. Therefore $\Gamma \vdash_\hickss \K_\agi \psi$ which implies that $\K_\agi \psi \in \Gamma$.
 \smallskip
     
   \setlength{\leftskip}{0cm}  \noindent \textsc{Case for $\C$.} Suppose $\varphi = \C \psi$. We distinguish the following cases depending on $\hickp$. \smallskip

   \setlength{\leftskip}{0.5cm} \noindent \textsc{Case for $\hickp \in \{\hick, \hickt\}$.} For the left-to-right direction suppose $\C \psi \in \Gamma$. Let $\Delta \in R^{(\Sigma, \hickp)^*}(\Gamma)$. Then there are prime theories $\Delta_0, \ldots, \Delta_n$ with $\Delta_0 = \Gamma$, $ \Delta_n = \Delta$ and for each $k \in [n]$ there exists $\agi_k \in \A$ with $\Delta_k \Rel_\agi^{(\Sigma, \hickp)} \Delta_{k+1}$. We prove that $\psi \in \Delta_n$ and $\C \psi \in \Delta_n$ by induction on $n$. For $n=0$, $\C \psi \in \Delta_0 = \Gamma$ by assumption. Moreover, Lemma~\ref{l: properties of prime theories}, Item~\ref{l: properties of prime theories, Item 3}. implies that $\varphi \in \Gamma$. For $n = k+1$ we have that $\psi \in \Delta_k$ and $\C \psi \in \Delta_k$ by induction hypothesis. By Lemma \ref{l: properties of prime theories}, Item \ref{l: properties of prime theories, Item 3}. $\K_\agi \C \psi \in \Delta_k$ for all $\agi \in \A$ and therefore $\K_{\agi_k} \C \psi \in \Delta_k$. Hence, $\C \psi \in \Delta_{k+1}$, implying that also $\psi \in \Delta_{k+1}$. Therefore $\psi \in \Delta$ for any $\Delta \in R^{(\Sigma, \hickp)^*}(\Gamma)$. By induction hypothesis $\fw^{(\Sigma, \hickp)}, \Delta \models \psi$ for any such $\Delta$ and thus $\fw^{(\Sigma, \hickp)}, \Gamma \models \C \psi$. 
     
     For the right-to-left direction suppose that $\fw^{(\Sigma, \hickp)}, \Gamma \models \C \psi$. For any $\Delta \in R^{(\Sigma, \hickp)^*}(\Gamma)$ thus holds that $\fw^{(\Sigma, \hickp)}, \Delta \models \psi$. By induction hypothesis $\psi \in \Delta$. In order to show that $\Gamma \vdash_\mathrm{P} \C \psi$, we claim that $\vdash_\mathrm{P} \gamma(\Gamma) \rightarrow \C \gamma(\Gamma)$. By the presence of the rule $\mathsf{Ind}$, it suffices to prove that $\vdash_\mathrm{P} \gamma(\Gamma) \rightarrow \E \gamma(\Gamma)$. We first show that $\vdash_\mathrm{P} \gamma(\Gamma) \rightarrow \K_\agi \gamma(\Gamma)$ for $\agi \in \A$. To that end note that for any $\Delta \in R^{(\Sigma, \hickp)^*}(\Gamma)$,
        \begin{equation*}
            \K_\agi^{-1} \Delta \vdash_\mathrm{P} \gamma(\Gamma).
        \end{equation*}
        In fact if $ \K_\agi^{-1} \Delta \not \vdash_\mathrm{P} \gamma(\Gamma)$ for some $\Delta \in R^{(\Sigma, \hickp)^*}(\Gamma)$, then by the Lindenbaum Lemma there exists a prime theory $\Omega$ such that $\K_\agi^{-1}\Delta \subseteq \Omega$ and $\Omega \not \vdash_\mathrm{P} \gamma(\Gamma)$. By construction $\Delta \Rel_\agi^{(\Sigma, \hickp)} \Omega$, implying that $\Omega \in R^{(\Sigma, \hickp)^*}(\Gamma)$. By Lemma \ref{l: properties characteristic formula IM}, $\Omega \vdash_\mathrm{P} \gamma(\Gamma)$ which is a contradiction. Thus $\K_\agi^{-1} \Delta \vdash_\mathrm{P} \gamma(\Gamma)$. By Corollary \ref{c: deduction theorem} and $\mathsf{Nec_\agi}$ we obtain
        \begin{equation*}
            \vdash_\mathrm{P} \K_\agi (\bigwedge \K_\agi^{-1} \Delta \rightarrow \gamma(\Gamma)).
        \end{equation*}
    Using the $\mathsf{K_\agi}$-axiom and $\mathsf{MP}$ yields
    \begin{equation*}
        \vdash_\mathrm{P} \K_\agi \bigwedge \K_\agi^{-1} \Delta \rightarrow \K_\agi \gamma(\Gamma).
    \end{equation*}
    By (\ref{e: truthlemma3}),
    \begin{equation*}
        \K_\agi \K_\agi^{-1} \Delta \vdash_\mathrm{P} \K_\agi \gamma(\Gamma).
    \end{equation*}
    Hence, by Lemma \ref{l: assumption weakening} we obtain $\Delta \vdash_\mathrm{P} \K_\agi \gamma(\Gamma)$. By Corollary \ref{c: deduction theorem} $\vdash_\mathrm{P} \chi(\Delta) \rightarrow \K_\agi \gamma(\Gamma)$. As $\Delta$ was arbitrary, we obtain $\vdash_\mathrm{P} \gamma(\Gamma) \rightarrow \K_\agi \gamma(\Gamma)$ by repeatedly using the (correct substitution instance of the) intuitionistic tautology $(p \rightarrow r) \rightarrow ((q \rightarrow r) \rightarrow ((p \vee q) \rightarrow r)$. Now since $\agi \in \A$ was arbitrary we obtain $\vdash_\mathrm{P} \gamma(\Gamma) \rightarrow \E \gamma(\Gamma)$ by repeatedly using the (correct substitution instance of the) intuitionistic tautology $(p \rightarrow q) \rightarrow ((p \rightarrow r) \rightarrow (p \rightarrow (q \wedge r)))$. Thus applying $\mathsf{Ind}$ to $\gamma(\Gamma) \rightarrow \E \gamma(\Gamma)$ yields $\vdash_\mathrm{P} \gamma(\Gamma) \rightarrow \C \gamma(\Gamma)$. \smallskip
    
    \noindent By Lemma \ref{l: properties characteristic formula IM}, $\Gamma \vdash_\mathrm{P} \gamma(\Gamma)$ and so we obtain
    \begin{equation*}
        \Gamma \vdash_\mathrm{P} \C \gamma(\Gamma).
    \end{equation*}
   Since $\psi \in \Delta$ for every $\Delta \in R^{(\Sigma, \hickp)^*}(\Gamma)$, Lemma \ref{l: properties characteristic formula IM} implies $\vdash_\mathrm{P} \gamma(\Gamma) \rightarrow \psi$. Applying $\mathsf{Mon}$ gives
   \begin{equation*}
       \vdash_\mathrm{P} \C \gamma(\Gamma) \rightarrow \C \psi.
   \end{equation*}
   Thus applying $\mathsf{MP}$ yields
   \begin{equation*}
       \Gamma \vdash_\mathrm{P} \C \psi
   \end{equation*}
   and hence $\C \psi \in \Gamma$. \smallskip

   \noindent \textsc{Case for $\hicks$.} The left-to-right direction is similar to the case for $\hickp \in \{\hick, \hickt\}$. The right-to-left direction is very similar, but uses the same construction as used in the case for $\K_\agi$, subcase $\hicks$, in this lemma. Suppose $\fw^{(\Sigma, \hicks)}, \Gamma \models \C \psi$. For any $\Delta \in R^{(\Sigma, \hicks)^*}(\Gamma)$ thus holds that $\fw^{(\Sigma, \hicks)}, \Delta \models \psi$. By induction hypothesis $\psi \in \Delta$. In order to show that $\Gamma \vdash_\hicks \C \psi$ we claim that $\vdash_\hicks \gamma(\Gamma) \rightarrow \C \gamma(\Gamma)$. By the presence of the rule $\mathsf{Ind}$ it suffices to prove that $\vdash_\hicks \gamma(\Gamma) \rightarrow \E \gamma(\Gamma)$. We first show that $\vdash_\hicks \gamma(\Gamma) \rightarrow \K_\agi \gamma(\Gamma)$ for $\agi \in \A$. To that end note that for any $\Delta \in R^{(\Sigma, \hicks)^*}(\Gamma)$,
 \begin{equation*}
     \K_\agi \K_\agi^{-1} \Delta \vdash_\hicks \gamma(\Gamma).
 \end{equation*}
 In fact if $\K_\agi \K_\agi^{-1} \Delta \not \vdash_\hicks \gamma(\Gamma)$ for some $\Delta \in R^{(\Sigma, \hicks)^*}(\Gamma)$, then by the Lindenbaum Lemma there exists a prime theory $\Omega$ such that $\K_\agi \K_\agi^{-1} \Delta \subseteq \Omega$ and $\Omega \not \vdash_\hicks \gamma(\Gamma)$. By construction $\Omega \in R^{(\Sigma, \hicks)^*}(\Gamma)$, contradicting Lemma \ref{l: properties characteristic formula IM}. Hence $ \K_\agi \K_\agi^{-1} \Delta \vdash_\hicks \gamma(\Gamma)$. Using the same argument as in the case for $\K_\agi$ starting from (\ref{e: completeness S4 K1}) we obtain
 \begin{equation*}\label{e: completeness S4 C1}
     \Delta \vdash_\hicks \K_\agi \gamma(\Gamma).
 \end{equation*}
 Corollary \ref{c: deduction theorem} yields $\vdash_\hicks \chi(\Delta) \rightarrow \K_\agi \gamma(\Gamma)$ and so $\vdash_\hicks \gamma(\Gamma) \rightarrow \K_\agi \gamma(\Gamma)$. We obtain $\vdash_\hicks \gamma(\Gamma) \rightarrow \E \gamma(\Gamma)$ by the same argument as in the previous case and so $\vdash_\hicks \gamma(\Gamma) \rightarrow \C \gamma(\Gamma)$ by $\mathsf{Ind}$. From this we deduce $\Gamma \vdash_\hicks \C \psi$ by the same argument as before. \smallskip

 \noindent \textsc{Case for $\hickss$.} The left-to-right direction is identical to the previous cases. For the right-to-left direction suppose $\fw^{(\Sigma, \hickss)}, \Gamma \models \C \psi$. For any $\Delta \in R^{(\Sigma, \hickss)^*}(\Gamma)$ thus holds that $\fw^{(\Sigma, \hickss)}, \Delta \models \psi$. By induction hypothesis $\psi \in \Delta$. In order to show that $\Gamma \vdash_\hickss \C \psi$ we claim that $\vdash_\hickss \gamma(\Gamma) \rightarrow \C \gamma(\Gamma)$. By the presence of the rule $\mathsf{Ind}$ it suffices to prove that $\vdash_\hickss \gamma(\Gamma) \rightarrow \E \gamma(\Gamma)$. We first show that $\vdash_\hickss \gamma(\Gamma) \rightarrow \K_\agi \gamma(\Gamma)$ for $\agi \in \A$. To that end note that for any $\Delta \in R^{(\Sigma, \hickss)^*}(\Gamma)$,
 \begin{equation*}
     \K_\agi \K_\agi^{-1} \Delta, \neg \K_\agi \neg \K_\agi^{-1} \vdash \gamma(\Gamma)
 \end{equation*}
 as otherwise there would exists a $\K$-maximally consistent prime theory $\Omega$ with $\K_\agi \K_\agi^{-1} \Delta \cup \neg \K_\agi \neg \K_\agi^{-1} \Delta \subseteq \Omega$ and $\Omega \not \vdash_\hickss \gamma(\Gamma)$. Note that $\Delta \Rel_\agi^{(\Sigma, \hickss)} \Omega$ and so we obtain that $\Omega \in R^{(\Sigma, \hickss)^*}(\Gamma)$, contradicting Lemma \ref{l: properties characteristic formula IM}. We obtain
 \begin{equation*}
     \Delta \vdash \K_\agi \gamma(\Gamma)
 \end{equation*}
 by the same arguments as used for the corresponding case starting from (\ref{e: completeness S5 K1}). From this we deduce $\Gamma \vdash_\hickss \C \psi$ as before.
\end{proof}

\end{document}
